\begin{proof}[Proof of Theorem~\ref{thm:crossing}]
The construction has three steps. We identify an integral pairing
of operators, factor it through states with two different energy
bounds, and count the states allowed by each bound.

\smallskip\noindent\emph{The pairing.}
Use the exterior algebra over $\mathbb Z$ on generators indexed by
$i\ge 0$. Its basis states are finite subsets of these indices,
written in increasing order. Let $c_i$ wedge the $i$th generator on
the left and let $a_i$ be its signed deletion. Creation operators
anticommute with one another, as do deletion operators, and
$a_i c_j+c_j a_i=\mathbf 1_{i=j}\mathrm{Id}$. Give a basis state $I$
energy $h(I)=\sum_{i\in I}(i+1/2)$. For background on the free fermion and vertex superalgebras see \cite{KacWang}. Introduce the fields
\[
 \psi(z)=\sum_{i\ge0}(c_i z^i+a_i z^{-i-1}),
 \qquad \psi_j(z)=\partial^{(j)}_z\psi(z),
 \qquad Y_I(z)={:}\prod_{i\in I}\psi_i(z){:}.
\]
Here $\partial^{(j)}_z z^b=\binom bj z^{b-j}$ for every integer $b$,
using the coefficient of $v^j$ in $(z+v)^b$. The product is ordered
by increasing $i$. The colons instruct us to expand and move creations
to the left of deletions using permutation signs only. Such a normal
product is skew under interchange of two fields. A coefficient $[z^b]Y_I$
shifts energy by $b+h(I)$. Indeed each summand of a differentiated
field shifts energy by its exponent plus the index of differentiation
plus $1/2$.

For equally large nonempty sets $I,J$, put $m=|I|=|J|$, $h=h(I)$,
$h'=h(J)$ and $L=h+h'-1$, a nonnegative integer. We claim
\begin{multline}\label{eq:crossing-field}
 (z-w)^L\big(Y_I(z)Y_J(w)-(-1)^{m^2}Y_J(w)Y_I(z)\big)
   \\= C_{IJ}\sum_{s\in\mathbb Z}z^{-s-1}w^s\,\mathrm{Id},
\end{multline}
where
\[
 C_{IJ}=(-1)^{\binom m2+\sum_{i\in I}i}
       \det\left[\binom{i+j}{i}\right]_{i\in I,\,j\in J}.
\]
To see this, normal order the product of the two blocks. When a
deletion from the first block meets a creation from the second, the
anticommutation relation permits either a signed interchange or a
scalar contraction. Summing the scalar for fields $i,j$ gives the
expansion in $w/z$ of
\[
  \partial^{(i)}_z\partial^{(j)}_w\frac{1}{z-w}
   =\frac{(-1)^i\binom{i+j}{i}}{(z-w)^{i+j+1}},
\]
since the undifferentiated contraction is
$\sum_{r\ge0}z^{-r-1}w^r$. Thus the expansion is indexed by partial
matchings between the blocks. Each summand is the product of the
matched scalars times the normal product of the unmatched fields.
Its sign is that of the permutation listing the matched pairs first,
in first-field/second-field order within each pair, followed by the
unmatched fields in their original order. In fact the sign for a
contraction is the sign of bringing the two operators together.
Moving the resulting pair as a block incurs no further sign.

For the reversed product, use the same ordering of pairs and the
same final order of the unmatched fields. Interchanging the original
blocks contributes $(-1)^{m^2}$. Reversing the orientation of any
pair contributes a further minus sign in the permutation, precisely
canceled by reversing its rational kernel, since
$(-1)^j/(w-z)^{i+j+1}=(-1)^{i+1}/(z-w)^{i+j+1}$. Consequently after
multiplication by $(-1)^{m^2}$ the reversed summands expand the same
rational kernels in $z/w$. For an incomplete matching an unmatched
field in each block accounts for at least one unit of missing pole
order. Its pole order is therefore at most $L$, and multiplication
by $(z-w)^L$ makes the two expansions equal. For a full matching the
pole order is $L+1$. Interleaving the two increasing lists contributes
$(-1)^{\binom m2}$, permuting partners contributes the determinant
sign, and the kernels contribute the remaining row signs in $C_{IJ}$.
The two expansions of $1/(z-w)$ differ by
$\sum_{s\in\mathbb Z}z^{-s-1}w^s$. This proves
\eqref{eq:crossing-field}.

We justify the formal products coefficientwise. Truncate the modes at
$N$ and first
perform the two normal-ordering expansions with their truncated scalar
kernels. Fix input and output states and a monomial $z^b w^{b'}$.
In the first product the intermediate energy is the input energy plus
$b'+h'$. In the reversed product the increment is $b+h$. There are
only finitely many states at either energy. Each individual normal
block between specified states uses deletion modes from its input
and creation modes from its output. Hence its matrix coefficient
stabilizes as a Laurent polynomial. The same argument applies to
the unmatched normal product, even in two variables. For each of its
finitely many monomials, prescribing the total $w$-power bounds all
contracted modes in the first expansion: a creation in mode $r$
differentiated $j$ times has $w$-power $r-j\ge0$ whenever its
coefficient is nonzero. The reversed expansion uses the $z$-power in
the same way. Both Wick expansions therefore pass coefficientwise to
the infinite-mode limit, where the rational-kernel cancellation above
applies. Multiplication by the fixed polynomial $(z-w)^L$ preserves
these limits. All coefficients are integers, including the divided
derivatives, so reduction modulo $2$ also preserves the identity.

\smallskip\noindent\emph{The factorizations.}
Let $D(H)$ count the states of energy at most $H$, with $D(H)=0$ for
$H<0$. Fix $H,E\ge1$ and an integer $s$. Use nonempty $I,J$ of
energy at most $H$, and input and output states $\beta,\alpha$ of
energy at most $E$. For a block of given $m,h,h'$ take
$[z^{-s-1}w^s]$ in \eqref{eq:crossing-field}. For the multiplier
monomial $z^{L-a}w^a$, $0\le a\le L$, the first product acts on
$\beta$ first by $[w^{s-a}]Y_J$, reaching energy at most
$E+h'+s-a\le E+H+s$. The reversed product acts first by
$[z^{-s-1-L+a}]Y_I$, reaching energy at most
$E-s-h'+a\le E+H-s$. Sum over intermediate states in each matrix
product. This writes $C_{IJ}\mathbf 1_{\alpha=\beta}$ as a sum of
two matrices, each factored over $\mathbb Z$. On the cut
$(I,\alpha)\mid(J,\beta)$ the first factors through at most
$(L+1)D(E+H+s)$ coordinates, indexed by $a$ and the intermediate
state. On the cut $(I,\beta)\mid(J,\alpha)$ the second factors
through at most $(L+1)D(E+H-s)$ coordinates. The integer coefficients
of the multiplier and the sign of the reversed product can be placed
on either factor.

Extend each block by zero and sum over $m,h,h'$. The number of blocks
and the largest $L+1$ are bounded by a polynomial in $H+1$, since
all energies are half-integers and $m\le 2H$. Define $C_{IJ}=0$ for
$|I|\ne |J|$. The resulting target is the full matrix
$C\otimes\mathrm{Id}_{D(E)}$ on the four indices. The two factorizations
have respective sizes at most
\[
  \operatorname{poly}(H+1)D(E+H+s),\qquad
  \operatorname{poly}(H+1)D(E+H-s).
\]
To turn $C$ into an identity, let $P_{i\ell}=\binom{i}{\ell}$
for nonnegative indices up to the largest mode under consideration.
Vandermonde's identity gives
$(PP^{\mathsf T})_{ij}=\sum_\ell\binom{i}{\ell}\binom{j}{\ell}
=\binom{i+j}{i}$. For fixed $m$, let $P^{[m]}$ be the matrix of
$m$-rowed minors restricted to sets of energy at most $H$ on both
indices. A minor of $P$ vanishes unless its increasing column list
is componentwise bounded by its increasing row list, since $P$ is
lower triangular. Thus restricting to this energy downset discards
no term of Cauchy--Binet, and the corresponding block of $C$ is
$P^{[m]}(P^{[m]})^{\mathsf T}$ up to row signs. In lexicographic
order $P^{[m]}$ is unit lower triangular. Every block is therefore
unimodular, and applying $C^{-1}$ at the $I$ index gives the
identity of dimension $D(H)-1$ without increasing either factorization
size.

\smallskip\noindent\emph{Counting the intermediate states.}
Write $\ln$ for the
natural logarithm. With
$A=\int_0^\infty\ln(1+e^{-x})\,dx>0$, the partition function is
\[
  \mathcal Z(\lambda)=\prod_{i\ge0}(1+e^{-\lambda(i+1/2)}),
  \qquad \ln\mathcal Z(\lambda)\sim A/\lambda
  \quad(\lambda\downarrow0).
\]
Under probabilities proportional to $e^{-\lambda h(I)}$ the occupation
variables are independent, with means $1/(e^{\lambda(i+1/2)}+1)$.
The mean energy is asymptotic to
$\lambda^{-2}\int_0^\infty x/(e^x+1)\,dx=A/\lambda^2$,
by integration by parts. Its variance is $O(\lambda^{-3})$ by
summing $(i+1/2)^2e^{-\lambda(i+1/2)}$. These Riemann-sum
estimates have exponentially decreasing tails. From
$D(H)\le e^{\lambda H}\mathcal Z(\lambda)$ at
$\lambda=\sqrt{A/H}$ we obtain the upper asymptotic bound. For the
lower bound take $0<\epsilon<1/2$ and
$\lambda=\sqrt{A/((1-\epsilon)H)}$. Chebyshev's inequality puts
probability $1-o(1)$ on energies between $(1-2\epsilon)H$ and $H$.
The number of states in this interval is at least
$(1-o(1))\mathcal Z(\lambda)e^{\lambda(1-2\epsilon)H}$.
Letting $\epsilon\downarrow0$ gives
\[
     \log_2 D(H)\sim c\sqrt H,\qquad c=2\sqrt A/\ln 2 .
\]
Finally take $H=(nt/c)^2$, $E=(nu/c)^2$ and
$s=\lfloor\delta n^2/c^2\rfloor$. Both shifted energy bounds are
positive for large $n$. Set $d_X=D(H)-1$, $d_Y=D(E)$ and identify
$(I,J,\alpha,\beta)$ with $(X_+,X_-,Y_+,Y_-)$. The polynomial
factors cost $o(n)$ in the logarithms, proving the theorem.
\end{proof}
